\section{Proof of the Cores Lemma}
\label{app:cores}

We prove \cref{lem:cores}. Let $w$ be nonempty with $m = \len{w}$, let $u_j$ and $\ell_j$ be its
hierarchy and markers, and let $\mathit{core}_j(x)$ be the core of position $x$ of $u_j$.

\paragraph{Part 1}
The cores of $u_{j+1}$ are a new start core followed by the cores of the nonfinal occurrences
of $\ell_j$ in $u_j$, in their order in $u_j$. By induction on $j$, the start cores of $u_j$
therefore form an initial segment of pairwise distinct cores, the ordinary cores increase from
left to right, and $\Cores_{j+1} \subseteq \Cores_j$. The start cores of $u_{j+1}$ are the new
one and those inherited from the start segment of $u_j$, so there are at most $1 + j$ of
them.

\paragraph{Shifted Expansions}
Let $\Gamma_0$ be the alphabet of $w$ and $\Gamma_{j+1} = (\Gamma_j \setminus \{\ell_j\})^*$, so that $u_j$ is a
word over $\Gamma_j$. Define $\xi_0(a) = a$ and
\[
  \xi_{j+1}([X]) = \xi_j(\ell_j)\,\xi_j(X[0]) \cdots \xi_j(X[b-1])
  \qquad\text{for } X = X[0]\cdots X[b-1].
\]
A set of nonempty words is a \emph{code} if every concatenation of its
elements has a unique factorization into those elements; prefix-freeness
is not required.
\begin{enumerate}[label=(\alph*),leftmargin=*]
\item \emph{$\xi_j$ is injective, and $\xi_j(\Gamma_j)$ is a code.}
  At level zero, the image consists of one-letter words. For the induction step, the map
  $[X] \mapsto \ell_j X$ has a code as its image, since each of its words starts with $\ell_j$
  and contains no other $\ell_j$. The extension of $\xi_j$ to words is an injective
  morphism because $\xi_j(\Gamma_j)$ is a code, and the image of a code under an injective
  morphism is a code.
\item \emph{Invariant.} Put $t_{-1} = m$. If $\Cores_j \ne \emptyset$, then every position
  after an ordinary position of $u_j$ is ordinary; for every ordinary position $x$ of $u_j$
  other than the last,
  $\wsub{w}{\mathit{core}_j(x)}{\mathit{core}_j(x+1)} = \xi_j(u_j[x])$; and for the last
  position, $\wsub{w}{t_j}{t_{j-1}} = \xi_j(\ell_j)$.

  The case $j = 0$ is immediate. For the step, let $x_0 < \cdots < x_{h-1}$ be the positions of
  $\ell_j$ in $u_j$, so that $x_{h-1}$ is its last position and position $i \ge 1$ of $u_{j+1}$
  holds the block formed by the positions $x_{i-1}+1, \dots, x_i-1$ and has the core
  $\mathit{core}_j(x_{i-1})$. If that core is ordinary, the positions $x_{i-1}, \dots, x_i - 1$
  of $u_j$ are ordinary and not last, so by induction the factor of $w$ from
  $\mathit{core}_j(x_{i-1})$ to $\mathit{core}_j(x_i)$ is
  $\xi_j(\ell_j)\,\xi_j(u_j[x_{i-1}+1]) \cdots \xi_j(u_j[x_i - 1]) =
  \xi_{j+1}(u_{j+1}[i])$. For $i < h-1$, the endpoint $\mathit{core}_j(x_i)$ is
  $\mathit{core}_{j+1}(i+1)$; for $i = h-1$, it is $\mathit{core}_j(x_{h-1}) = t_j$, and the
  letter $u_{j+1}[h-1]$ is the marker $\ell_{j+1}$.
\end{enumerate}
Write $r_j = \wsuf{w}{t_j}$ for the terminal core word of level $j$ and $r_{-1} = \eps$. By
(b), $r_j = \xi_j(\ell_j)\, r_{j-1}$.

\paragraph{Part 2}
We show by induction on $j$ that, if $\Cores_j \ne \emptyset$, then
$\Cores_{j+1} = \{s \mid s < t_j,\ r_j \pre \wsuf{w}{s}\}$. (If $\Cores_j = \emptyset$, then
$\Cores_{j+1} = \emptyset$ by Part~1.)

($\subseteq$) Let $s \in \Cores_{j+1}$. Then $s$ is the core of a nonfinal occurrence $x$ of
$\ell_j$ in $u_j$, so $s < t_j$, and by (b) the factor of $w$ from $s$ to
$\mathit{core}_j(x+1)$ is $\xi_j(\ell_j)$. For $j = 0$ this factor is $r_0$. For $j > 0$, the
core $\mathit{core}_j(x+1)$ belongs to $\Cores_j$ and hence, by the induction hypothesis, is an
occurrence of $r_{j-1}$. So $\wsuf{w}{s}$ starts with $\xi_j(\ell_j)\,r_{j-1} = r_j$.

($\supseteq$) Let $s < t_j$ with $r_j \pre \wsuf{w}{s}$. For $j = 0$, $r_0 = \ell_0$ is the last
letter of $w$, and $s$ is a nonfinal occurrence of it, so $s \in \Cores_1$. Let $j > 0$. By the
induction hypothesis, the occurrences of $r_{j-1}$ in $w$ are exactly the elements of
$\Cores_j \cup \{t_{j-1}\}$ (an occurrence $i$ satisfies $i \le t_{j-1}$, since
$\len{r_{j-1}} = m - t_{j-1}$). As $t_j \in \Cores_j$, the word $r_{j-1}$ is a prefix of $r_j$,
so $s$ is an occurrence of $r_{j-1}$ with $s < t_{j-1}$; thus $s \in \Cores_j$, say
$s = \mathit{core}_j(x)$, and $x$ is not the last position because $s < t_j$. Inside the
suffix $r_j$ of $w$, the occurrences of $r_{j-1}$ are at $t_j$ and $t_{j-1}$ only, because
$t_j = \max \Cores_j$. Hence the word $r_j$ contains $r_{j-1}$ exactly at the offsets $0$ and
$\varpi = t_{j-1} - t_j$. Consequently the occurrences of $r_{j-1}$ in $w$ that start in
$[s, s+\varpi]$ are exactly $s$ and $s+\varpi$: an occurrence starting strictly between them
ends before $s + \len{r_j}$ and would be a third occurrence inside the copy of $r_j$ at~$s$.
Since $s + \varpi < t_{j-1}$, the next element of $\Cores_j$ after $s$ is $s + \varpi$, and by
Part~1 it is $\mathit{core}_j(x+1)$. By (b),
$\xi_j(u_j[x]) = \wsub{w}{s}{s+\varpi} = \wsub{w}{t_j}{t_{j-1}} = \xi_j(\ell_j)$, and
injectivity gives $u_j[x] = \ell_j$. Since $x$ is not the last position, $s \in \Cores_{j+1}$.
\qed

\section{Details of the Bounded-Depth Compression}
\label{app:search}

We prove the dictionary normal form and the compression statements of \cref{sec:search}
for an arbitrary lettered monitor with $N\ge1$ states. Its nonempty holes
supply every marker used below; no finiteness assumption on the alphabet
is needed. All dictionaries and tables considered here are finite.

\subsection{Comparisons and Isolated Exchange}

\begin{lemma}[Comparison-preserving maps]\label{lem:comparison-map}
Let $v$ be a nonempty word and $\theta$ a letter map. Suppose that $\theta$ identifies
no unequal predecessor pair of any comparison, including terminal ones:
\[
 1\le s<e\le|v|,\quad v[e:]\pre v[s:],\quad
 \theta(v[s-1])=\theta(v[e-1])\quad\Longrightarrow\quad v[s-1]=v[e-1].
\]
Then $v$ and $\theta(v)$ have exactly the same comparisons. Moreover, on the
letters of $v$, $\theta(a)=\theta(\ell)$ implies $a=\ell$, where $\ell$ is the
last letter of $v$.
\end{lemma}
\begin{proof}
A letter map preserves every old comparison. If $(s,e)$ is a new one,
choose the largest offset $t<|v|-e$ at which
$v[s+t]\ne v[e+t]$. The remaining suffixes agree, so
$(s+t+1,e+t+1)$ is an old comparison. Its unequal predecessors have equal
images, contrary to the hypothesis. To prove the last assertion, an
occurrence of $a$ before the final position gives a terminal comparison
between the cuts following that occurrence and the last letter. Its
predecessors cannot be identified unless they agree. The final occurrence
itself is immediate.
\end{proof}

We first justify the state information used by the exchange types.
The state set at level $i$ is $Q$ with $i$ successive symbolic states
adjoined. Ordinary states form a closed copy of $Q$: at level zero this
is immediate, and a derived transition from an ordinary state runs the
previous monitor over $\ell X$, hence stays ordinary by induction.
The shifted-run definition also identifies the entry state of a position
with its core: an ordinary core $p$ carries the original state at cut $p$;
a symbolic core carries the corresponding adjoined state. Indeed, the
first block starts at the new symbolic state, and each later block starts
at the entry state of its preceding marker. This is precisely the core
recursion of \cref{lem:cores}. Thus all occurrences of an unprotected
letter have ordinary entry states, and every entry state occurring in
$u_i$ belongs to $Z_i$.

\begin{proof}[Proof of \cref{lem:exchange-types}]
The transition table has $N^N$ possibilities. The target table uses $N$
bits, the two self tables $2N^2$ bits, the protected tables
$4N\,|Z_i|\,|\Pi_i|$ bits, and the initial rows $2N$ bits. Hence the number
of types is at most
\[
 N^N2^{2N^2+4N|Z_i||\Pi_i|+3N}.
\]
Using $|\Pi_i|\le s_i+1\le d+1$ and $|Z_i|\le N+s_i\le N+d$ gives the bound
$\TyBound(N,d)$ of \eqref{eq:dictionary-constants}.
\end{proof}

For the local part of \cref{lem:isolated-exchange}, fix $i$, put $v=u_i$,
and let $\theta$ identify $a$ with $a'$ and fix every other letter.
An unequal predecessor pair identified by $\theta$ would be $(a,a')$ or
$(a',a)$. A nonterminal such pair contradicts isolation; a terminal one
would contain the protected marker. \Cref{lem:comparison-map} therefore
preserves comparisons. Induction on positions preserves the full state
run: the only changed transitions read $a$ at ordinary states, where its
transition table equals that of $a'$. The final letter is protected and
unchanged, so its entry state and target test are unchanged.

For each comparison, consider its predecessor labels. If neither letter
is $a$, the labels are unchanged. If both are $a$, the relevant self table
preserves admission. If exactly one is $a$, the other is protected:
for a nonterminal comparison this follows from isolation, and for a
terminal comparison its later predecessor is the marker. The appropriate
protected row, in the appropriate orientation, preserves either admission
predicate. A comparison starting at zero uses the initial-label row;
at positive levels this row is false. Protected letters at symbolic
states are covered by the $Z_i$ columns. These cases cover all comparisons,
so $\theta(v)$ is a hole of the same level-$i$ monitor. Expanding it with the
markers of the original hierarchy below level $i$ gives a nonempty base hole
by \cref{lem:der-holes}.
It remains to prove the asserted structural properties of that expansion.

\subsection{Induced Level Maps}

Continue a hierarchy past its first unary level when discussing all
levels; desubstitution preserves unary words. For a level map $\theta$ define
$\theta_0=\theta$ and $\theta_{r+1}(X)=\theta_r^*(X)$, where the star denotes letterwise
application to a block body.

\begin{lemma}[Shape under a level map]\label{lem:level-maps}
Let $w$ be a nonempty hole, and let $\theta$ map the dictionary of $u_i$ into
itself. Suppose that it identifies no unequal predecessor pair, as in
the hypothesis of \cref{lem:comparison-map}, and makes $\theta(u_i)$ a hole
of $\mathcal M_i$. Expand $\theta(u_i)$ with the
markers of the original hierarchy below level $i$ to obtain $w'$. Then $w'$
is a hole, its markers below $i$ are unchanged, and its dictionaries below
$i$ are subsets of the old ones. At every $r\ge0$,
\[
 u'_{i+r}=\theta_r(u_{i+r}),\qquad
 \ell'_{i+r}=\theta_r(\ell_{i+r}),
\]
and $\theta_r$ identifies no unequal predecessor pair of a comparison of
$u_{i+r}$. Exactly the same levels are unary. In particular, depth is
preserved and no dictionary grows.
\end{lemma}

\begin{proof}
Every image letter occurs in the old dictionary, so it has a valid
expansion avoiding each lower marker. Expansion and cutting are inverse
at every lower level: a block body avoids its marker and its following
marker is present. Thus the decoded hierarchy reaches $\theta(u_i)$ exactly,
with unchanged lower markers. Repeated use of \cref{lem:der-holes} proves
holehood. Descending through these expansions also proves dictionary
containment below $i$.

We verify the upward induction. For any map $F$ satisfying
the hypothesis of \cref{lem:comparison-map} on a word $u$ ending in $\ell$, \cref{lem:comparison-map} reflects the
marker: no other letter maps to $F(\ell)$. Hence cutting at markers
commutes with $F$. Write $V$ for the block word of $u$. The induced
map $F^*$ on its letters identifies no unequal predecessor pair of an
upper comparison. To see this, suppose the cuts following blocks $X,Y$
form a comparison of $V$, and $F^*(X)=F^*(Y)$. In the image of $u$,
the cuts at the beginnings of these blocks form a comparison: their
first blocks agree and the following block suffixes compare.
Comparison preservation by $F$ gives the same lower comparison before
mapping. Since $X,Y$ avoid $\ell$, its first aligned markers force
$X=Y$. This proves the required nonidentification property upstairs,
including when the upper comparison is terminal. Induction gives the
claimed level words, markers, and comparison sets at every higher level.

Finally, marker reflection implies that a higher mapped word is unary
iff the old word is unary: if every image letter equals the image marker,
every old letter is the marker. If $u_i$ is nonunary, no lower word in
either hierarchy is unary, since a unary word has only unary derived
levels. If $u_i$ is unary, its dictionary is a singleton and $\theta$ maps
it into itself, so decoding makes $w'=w$. This proves the claim at
lower levels as well. Higher dictionaries are images of old dictionaries,
so their sizes cannot increase.
\end{proof}

\begin{proof}[Completion of \cref{lem:isolated-exchange}]
The local argument already proved comparison preservation and holehood
at level $i$. The map sends the dictionary into itself, because $a'$ occurs
in $u_i$. Apply \cref{lem:level-maps}. The new level-$i$ dictionary is the old
one with $a$ removed, so its size strictly decreases. Write $|a|_0$ and
$|a'|_0$ for the lengths of the ordinary level-zero expansions of $a$ and
$a'$ with the markers of the original hierarchy. The decoded length changes by
\[
 \#_a(u_i)(|a'|_0-|a|_0).
\]
Both merge directions satisfy the hypotheses. Choose the direction with
$|a'|_0\le|a|_0$ to make this difference nonpositive.
\end{proof}

The level maps establish hierarchy geometry, not agreement of all the
higher derived predicates under noninjective maps. After decoding, derive
the monitors again from the new word. Holehood in those monitors follows
from \cref{lem:der-holes}. Protected sets, isolation, and exchange types
are likewise those of the new word.

\subsection{Dictionary Normal Form, Boundary Pairs, and the Recurrence}

\begin{proof}[Proof of \cref{lem:dictionary-normal-form}]
Use the natural-number potential $\sum_{j=0}^{\Delta}g_j$, the total
size of the dictionaries. Among the holes of the same depth $\Delta$ as the
given hole and of no greater length, choose one of least potential. It is
dictionary-normal. Otherwise, at some level $1\le i<\Delta$ of its own
hierarchy, two distinct isolated letters have the same exchange type; merge
them in the direction of smaller expansion. By \cref{lem:isolated-exchange},
the result is a hole of depth $\Delta$ and no greater length, no dictionary
grows, and one strictly shrinks. Its potential is therefore strictly
smaller, contradicting the choice. The hierarchy, the protected sets, and
the exchange types in this argument are those of the chosen hole, so
dictionary normality holds at every positive nonunary level simultaneously.
\end{proof}

\begin{proof}[Counting part of \cref{lem:dictionary-count}]
Partition $\mathcal A_i$ into protected letters, isolated unprotected
letters, and the remaining letters. The first part has size at most
$d+1$. Dictionary normality injects the second part into the type space, so it has
size at most $\TyBound(N,d)$. Each letter in the third part is a component of
some ordered boundary pair with both components unprotected. The union
of the two projections of $\operatorname{BP}_i$ has size at most $2b_i$.
Adding gives $g_i\le \AlBound(N,d)+2b_i$.
\end{proof}

\begin{proof}[Propagation part of \cref{lem:dictionary-count}]
For distinct blocks $X,Y$ avoiding a marker $\ell$, define
$\mu_\ell(X,Y)$ as their first unequal letters when scanned from right
to left, using $\ell$ when one body is exhausted. Equivalently, remove
their longest common suffix; take the last remaining letter in each body,
or $\ell$ for an empty remainder. These two letters are unequal.
Set $\mu_\ell(X,X)=(\ell,\ell)$ to make the map total.

The \emph{endpoint pairs} of level $i$ are the predecessor pairs of the
left-maximal comparisons that end at the final cut:
\[
 \operatorname{EP}_i=\{(u_i[s-1],u_i[e-1]):
       (s,e)\text{ is left maximal},\ e=|u_i|\}.
\]
Their second component is the marker $\ell_i$.

Take a lower boundary pair witnessed by $(s,e)$ in the expansion of a
block word $V$, with $e$ nonterminal. By \cref{lem:alignment}, the cuts
lie in blocks $X=V[i_s]$ and $Y=V[i_e]$, with $i_s<i_e$, the same
remaining body suffix, and an upper comparison $(i_s+1,i_e+1)$.
Because $s\ge1$, the predecessor of a cut at a block's beginning is the
preceding marker; a cut inside a body has its ordinary predecessor.
The unequal lower predecessors are exactly $\mu_\ell(X,Y)$: the common
suffix extends right to the marker and cannot extend one more letter left.
In particular $X\ne Y$. Thus $(X,Y)$ belongs either to the upper
nonterminal boundary pairs or to its endpoint pairs, and
\[
 \operatorname{BP}_i\subseteq
 \mu_{\ell_i}(\operatorname{BP}_{i+1}\cup\operatorname{EP}_{i+1}).
\]
This also covers lower pairs with one component equal to $\ell_i$ and
upper comparisons ending at the final cut. Every upper endpoint pair
has the form $(a,\ell_{i+1})$ with $a\ne\ell_{i+1}$. There are at most
$g_{i+1}-1$ such pairs. Taking cardinalities proves
\eqref{eq:boundary-propagation}.
\end{proof}

\begin{proof}[Proof of \cref{lem:dictionary-bound}]
All counts now refer to one fixed dictionary-normal hole; abbreviate
$\AlBound=\AlBound(N,d)$. Its unary top has
$g_\Delta=1$ and $b_\Delta=0$. If $\Delta\ge1$, propagation gives $b_{\Delta-1}=0$.
For $1\le i<\Delta-1$, counting at level $i+1$ and propagation give
\[
 b_i\le b_{i+1}+g_{i+1}-1\le3b_{i+1}+\AlBound-1.
\]
Induct downwards from $\Delta-1$. If $r=\Delta-1-i$ and the bound holds at $i+1$,
\[
 2b_i\le3(\AlBound-1)(3^{r-1}-1)+2(\AlBound-1)=(\AlBound-1)(3^r-1).
\]
The base case has both sides zero. Adding $\AlBound$ gives
$g_i\le \AlBound+2b_i\le1+(\AlBound-1)3^r$, as claimed.
Since $r\le d$ and $\AlBound\ge1$, this is at most $\DictBound(N,d)$.
The unary dictionary has size one, also at most $\DictBound(N,d)$.
Put an arbitrary hole in dictionary normal form using \cref{lem:dictionary-normal-form}
to obtain the existential assertion. For $\Delta=0$ there are no positive
levels to bound; for $\Delta=1$ only the unary dictionary is needed.
\end{proof}

\subsection{Unary and Column Shortening}
Fix a lettered monitor $\mathcal M$ with state set $\mathcal Q$ and
start state $\zeta_0$, and a marker $\ell$. Write
$\mathit{lab}_0(\zeta)$ for the initial shifted-block label:
$(\zeta,\ell)$ for any $\zeta\in\mathcal Q$, and
$\triangleright$ at the fresh entry symbol $\star$.

\paragraph{Unary Words}
Let $v = z^L$ be a hole of $\mathcal{M}$, with states $\zeta_0, \zeta_1, \dots$ where
$\zeta_{p+1} = \delta(\zeta_p,z)$. Since every pair of cuts of a unary word is a comparison, $v$ is a
hole iff $\Tfin(\zeta_{L-1},z)$ holds, no label of a cut $s < L$ is related by $\Lfin$ to the
last label, and no two labels of cuts $s < e < L$ are related by $\Lambda$. If
$\zeta_i = \zeta_{i'}$ for some $i < i' \le L-1$, delete the letters between cuts $i$ and $i'$.
The cut $p > i$ of the new word has the state and label of the cut $p + (i'-i)$ of $v$, so the
labels of the new word form a subsequence of those of $v$ with the same last label. Every
condition above is therefore preserved. Repeating, we obtain a unary hole of length at most
$\len{\mathcal Q}$.

\paragraph{The Column Automaton}
Let $X_1,\dots,X_h$ be blocks avoiding $\ell$. The summary of \cref{sec:search} is a tuple of
four functions, where $\mathbb{B}$ is the set of Boolean values:
\begin{align*}
  \mathit{exit}_a(\zeta) &\in \mathcal Q, &&\text{the exit state of the shifted block of } X_a
                                 \text{ read from } \zeta,\\
  \mathit{term}_a(\zeta,\zeta'') &\in \mathbb{B}, &&\text{whether some label of it is related by }
                                 \Lfin \text{ to } (\zeta'',\ell),\\
  \mathit{eq}(a,b) &\in \mathbb{B}, &&\text{whether } X_a = X_b,\\
  \mathit{adm}(a,b,\zeta,\zeta') &\in \mathbb{B}, &&\text{whether } \Lambda_\ell \text{ relates }
                                 (\zeta,X_a) \text{ and } (\zeta',X_b),
\end{align*}
for blocks $a,b$, states $\zeta,\zeta' \in \mathcal Q \sqcup \{\star\}$, and states $\zeta'' \in \mathcal Q$. Its number
of values is the function $\mathit{col}(\len{\mathcal Q},h)$ of \eqref{eq:col}.

Right-align the blocks, pad them on the left, and read the resulting word of columns, each an
$h$-tuple of optional letters, from left to right (\cref{fig:columns}). The tuple of $h$ empty blocks has
the initial summary with $\mathit{exit}_a(\zeta) = \delta(\zeta,\ell)$ for $\zeta \in \mathcal Q$ and
$\mathit{exit}_a(\star) = \zeta_0$, with $\mathit{term}_a(\zeta,\zeta'') =
\Lfin(\mathit{lab}_0(\zeta),(\zeta'',\ell))$, with all blocks equal, and with
$\mathit{adm}(a,b,\zeta,\zeta') = \Lambda(\mathit{lab}_0(\zeta),\mathit{lab}_0(\zeta'))$. Reading a column
$\gamma$ updates all components simultaneously, using the old values on the right-hand sides. If $\gamma_a$ is absent, block $a$ has not started and its
components are unchanged. If $\gamma_a = z$, then
\[
  \mathit{exit}'_a(\zeta) = \delta(\mathit{exit}_a(\zeta),z),\qquad
  \mathit{term}'_a(\zeta,\zeta'') = \mathit{term}_a(\zeta,\zeta'') \vee
     \Lfin\bigl((\mathit{exit}_a(\zeta),z),(\zeta'',\ell)\bigr).
\]
Equality is updated to $\mathit{eq}(a,b) \wedge \gamma_a = \gamma_b$. For the admission bit, if
$\gamma_a = z$ and $\gamma_b = z'$ are both present,
\[
  \mathit{adm}'(a,b,\zeta,\zeta') =
  \Lambda\bigl((\mathit{exit}_a(\zeta),z),(\mathit{exit}_b(\zeta'),z')\bigr)
  \;\vee\; \bigl(z = z' \wedge \mathit{adm}(a,b,\zeta,\zeta')\bigr).
\]
If only $\gamma_a=z$ is present, set
\[
 \mathit{adm}'(a,b,\zeta,\zeta')=
 \Lambda\bigl((\mathit{exit}_a(\zeta),z),\mathit{lab}_0(\zeta')\bigr).
\]
If only $\gamma_b=z'$ is present, set
\[
 \mathit{adm}'(a,b,\zeta,\zeta')=
 \Lambda\bigl(\mathit{lab}_0(\zeta),(\mathit{exit}_b(\zeta'),z')\bigr).
\]
In these cases the absent track is still empty, so only the empty common
suffix is possible and no old admission bit is retained. If both entries
are absent, retain the old admission bit. The update is correct because $\Lambda_\ell$ compares positions at the \emph{same
distance from the ends} of the two blocks, with equal letters after them: appending a column
either relates the two new last positions, or extends the common suffix of an earlier related
pair by the letter $z = z'$. This is why the blocks are aligned at their right ends.

\paragraph{Shortening}
A word of columns is \emph{valid} if each track consists of absent entries followed by
letters. Deleting the columns between two cuts with the same summary preserves validity,
since a subsequence of a valid word is valid, and preserves the final summary. Hence every
tuple of blocks can be replaced by one with the same summary and all lengths below
$\mathit{col}(\len{\mathcal Q},h)$.

The summary determines the derived transition $\delta_\ell$, the target predicate
$\Tfin_\ell$, both derived relations on the labels with these letters, and the equalities
between the letters; that is, it determines the restriction of $\mathcal{M}_\ell$ to the letters
in use up to renaming. Replacing the blocks
therefore maps holes of $\mathcal{M}_\ell$ to holes. Since the replacement is injective on the
letters in use and preserves marker avoidance, it keeps a unary level of the canonical
hierarchy unary.


\subsection{Agreement Under Derivation}
For lettered monitors $\mathcal U,\mathcal V$ with the same state set,
a finite letter support $\mathcal E$, and a map $\rho$ from $\mathcal E$ into the alphabet
of $\mathcal V$, write $\mathcal U\equiv_{\mathcal E,\rho}\mathcal V$ for
\emph{agreement}: their start states coincide, $\rho$ is injective on
$\mathcal E$, and, for every $a\in \mathcal E$ and every state $\zeta$,
\begin{equation}\label{eq:monitor-agreement}
\begin{aligned}
\delta_{\mathcal V}(\zeta,\rho(a))&=\delta_{\mathcal U}(\zeta,a),\\
\Tfin_{\mathcal V}(\zeta,\rho(a))&\iff\Tfin_{\mathcal U}(\zeta,a),\\
\Lambda_{\mathcal V}(\widehat\rho L,\widehat\rho L')&\iff
  \Lambda_{\mathcal U}(L,L'),\\
\Lfin_{\mathcal V}(\widehat\rho L,\widehat\rho L')&\iff
  \Lfin_{\mathcal U}(L,L').
\end{aligned}
\end{equation}
The last two lines range over all labels whose letters lie in $\mathcal E$, as
well as the start label. Here $\widehat\rho(\triangleright)=\triangleright$
and $\widehat\rho(\zeta,a)=(\zeta,\rho(a))$; states are not renamed.
Agreement preserves runs, labels up to $\widehat\rho$, and holes of words
over $\mathcal E$. Indeed, induction gives the same state at every cut; injectivity
preserves both equality and inequality of letters, hence every comparison;
and \eqref{eq:monitor-agreement} preserves the target and all admissions.
Agreements compose: if $\mathcal U\equiv_{\mathcal E,\rho}\mathcal V$ and
$\mathcal V\equiv_{\rho(\mathcal E),\rho'}\mathcal W$, then
$\mathcal U\equiv_{\mathcal E,\rho'\circ\rho}\mathcal W$.

\begin{lemma}[Agreement under derivation]\label{lem:agreement-derivative}
Suppose $\mathcal U\equiv_{\mathcal E,\rho}\mathcal V$ and $\ell\in \mathcal E$.
For any finite family $F$ of blocks over $\mathcal E$ avoiding $\ell$, the
letterwise map $\rho^*$ gives
\[
 \mathcal U_\ell\equiv_{F,\rho^*}\mathcal V_{\rho(\ell)}.
\]
The image blocks avoid $\rho(\ell)$.
\end{lemma}
\begin{proof}
Injectivity on $\mathcal E$ implies injectivity on words over $\mathcal E$ and preserves
marker avoidance. The shifted runs start in the same state: for an
ordinary entry state, use agreement on the marker before the block;
for $\star$, use equality of start states. Induction on the block gives
equal exits and labels related by $\widehat\rho$ at every position.
A common-suffix decomposition maps to one after renaming. Conversely,
any common suffix of two renamed blocks pulls back to equal suffixes,
by injectivity. Thus both directions of the derived $\Lambda$ agree.
The terminal-comparison test ranges over the same positions and the same
exit states, so its truth is preserved by the last line of
\eqref{eq:monitor-agreement}. This proves preservation of both the
derived target and terminal relation. Relations involving the derivative
start label are false in both monitors.
\end{proof}

Fix a base monitor $\mathcal M_0$ with $N$ states and markers
$\boldsymbol\ell=(\ell_0,\ldots,\ell_{j-1})$. They determine a tower
$\mathcal M^{\boldsymbol\ell}_0=\mathcal M_0$ and
$\mathcal M^{\boldsymbol\ell}_{i+1}
=(\mathcal M^{\boldsymbol\ell}_i)_{\ell_i}$.
A finite family of level-$j$ letters is \emph{valid for the markers}
if every block avoids its level's marker and its ordinary expansion,
recursively, uses valid lower-level letters. Write
$\operatorname{Exp}_{\boldsymbol\ell,j}(a)$ for the level-zero word
obtained by expanding a single level-$j$ letter, with each marker
placed after its block. At level zero this is the one-letter word $a$.

\begin{lemma}[Supported simultaneous compression]\label{lem:simultaneous-compression}
Fix a nonempty word $w$, its original hierarchy and markers
$\boldsymbol\ell$, and $j\le d$. Suppose $|\mathcal A_i|\le g$ for
$1\le i\le j$. Put $K=\mathit{col}(N+d+1,g)$. For a finite set
$\mathcal E_j\subseteq\mathcal A_j$, there are new markers
$\widetilde{\boldsymbol\ell}$ and an injective renaming $\rho_j$ on $\mathcal E_j$
such that
\[
 \mathcal M^{\boldsymbol\ell}_j
 \equiv_{\mathcal E_j,\rho_j}\mathcal M^{\widetilde{\boldsymbol\ell}}_j,
 \qquad
 |\operatorname{Exp}_{\widetilde{\boldsymbol\ell},j}(\rho_j(a))|
 \le K^j\quad(a\in \mathcal E_j).
\]
The image family is valid for the new markers. Every recursive support
is a set contained in its original dictionary; at level zero there is no
cardinality assumption. Agreement is at every state, not only states
visited by the fixed word.
\end{lemma}
\begin{proof}
Induct on $j$, keeping $w$, its dictionaries, and $d,g$ fixed. At zero,
take the identity map and old markers. A letter has expansion size
$1=K^0$, regardless of the size of $\mathcal E_0$.

For the step from $j$ to $j+1$, the set $\mathcal E_{j+1}$ has at most $g$ elements.
Let $K_j=\mathit{col}(N+j,g)$ and
$\mathcal U=\mathcal M^{\boldsymbol\ell}_j$.
The column construction, with empty padding tracks if necessary,
replaces every $X\in \mathcal E_{j+1}$ by a subsequence $\tau(X)$ of length less
than $K_j$. Equality bits make $\tau$ injective on $\mathcal E_{j+1}$; the other fields
preserve transitions, target, and both admissions at every entry state.
Thus
\[
 \mathcal U_{\ell_j}\equiv_{\mathcal E_{j+1},\tau}\mathcal U_{\ell_j}.
\]
Form the set
\[
 \mathcal E_j = \{\ell_j\}\cup\bigcup_{X\in \mathcal E_{j+1}}\operatorname{alph}(\tau(X)).
\]
The marker belongs to $\mathcal A_j$. Each $X\in\mathcal A_{j+1}$
is an original block body of $u_j$, so all its letters belong to
$\mathcal A_j$. Since $\tau(X)$ is a subsequence of $X$, this proves
$\mathcal E_j\subseteq\mathcal A_j$. Counting distinct letters ensures
$|\mathcal E_j|\le g$ whenever $j\ge1$. At $j=0$ the
identity base case applies without such a bound.

Apply the induction hypothesis to $\mathcal E_j$, obtaining lower markers,
a monitor $\mathcal V=\mathcal M^{\widetilde{\boldsymbol\ell}}_j$,
and $\rho_j$ with agreement on $\mathcal E_j$ and expansion bound $K^j$.
Set
\[
 \widetilde\ell_j=\rho_j(\ell_j),\qquad
 \rho_{j+1}(X)=\rho_j^*(\tau(X)).
\]
Changing marker $j$ does not affect the monitors through level $j$.
By \cref{lem:agreement-derivative}, lower agreement lifts to agreement
on the shortened blocks between $\mathcal U_{\ell_j}$ and
$\mathcal V_{\widetilde\ell_j}$. Composition with column agreement
proves the displayed invariant at level $j+1$.

Both the marker and all body letters lie in $\mathcal E_j$. Injectivity therefore
preserves marker avoidance; validity at smaller levels follows from the
induction hypothesis. Each body followed by its marker has at most
$K_j$ letters. Formula \eqref{eq:col} is increasing in its positive state
argument, so $K_j\le K$. Expanding these letters gives at most
$K_jK^j\le K^{j+1}$ base letters. This proves agreement, validity,
containment, and the expansion bound together. No dictionary or normality
property of an intermediate compressed word was used.
\end{proof}

\begin{proof}[Proof of \cref{thm:search}]
Convert the reader to its lettered monitor, or start with any lettered
monitor. Put the given nonempty hole in dictionary normal form. By
\cref{lem:dictionary-bound}, its positive dictionaries through its first
unary level $\Delta\le d$ are bounded by $g=\DictBound(N,d)$. Keep this word
fixed as the original hierarchy for \cref{lem:simultaneous-compression}, and put
$K=\mathit{col}(N+d+1,g)$.
The unary argument gives a hole $z^L$ of its level-$\Delta$ monitor with
$1\le L\le N+\Delta$. Compress $\mathcal E_\Delta=\{z\}$ using the lemma. Agreement makes
$\rho_\Delta(z)^L$ a hole of the new level-$\Delta$ monitor. Validity allows exact
expansion by \cref{lem:der-holes}, yielding a nonempty base hole.
Cutting that expansion recovers the renamed unary word at level $\Delta$,
so its depth is at most $\Delta$. Its length is bounded by
\[
 LK^\Delta\le(N+\Delta)K^\Delta\le(N+d+1)K^d=\holeBound(N,d).
\]
Here $K\ge1$ since $N\ge1$. This also proves the depth-zero case,
where the expansion is the identity.
\end{proof}

\subsection{Closed-Form Estimate}
\begin{lemma}[Elementary estimate]\label{lem:elementary-estimate}
Put
\[
 \Xi(N,d)=3N^2+4N(N+d)(d+1)+3N+3d+4.
\]
Then $\DictBound(N,d)\le2^{\Xi(N,d)}$ and
$\holeBound(N,d)\le2^{2^{5\Xi(N,d)+3}}$.
For $d_N=2N(4N-1)-1$ and $N\ge1$,
\begin{equation}\label{eq:elementary-closed}
 \holeBound(N,d_N)\le2^{2^{1280(N+1)^5}}.
\end{equation}
\end{lemma}
\begin{proof}
Write $\Xi=\Xi(N,d)$ and
$\Xi_0=3N^2+4N(N+d)(d+1)+3N$.
The inequality $N\le2^N$ gives
$\TyBound(N,d)\le2^{\Xi_0}$. For nonnegative $p,q$, the sum of quantities
bounded by $2^p,2^q$ is at most $2^{p+q+1}$. Since
$d+1\le2^{d+1}$, we obtain $\AlBound(N,d)\le2^{\Xi_0+d+2}$.
Using $3^d\le2^{2d}$ gives
$\DictBound(N,d)=2+\AlBound(N,d)3^d\le2^{\Xi_0+3d+3}\le2^\Xi$.

Put $M=N+d+1$, $g=\DictBound(N,d)$, and $P=2^\Xi$. We have $M+1\le \Xi$,
so $M,M+1,d,g\le P$. Replacing $M$ by $2^M$ in the first factor
of \eqref{eq:col} yields
\[
 \mathit{col}(M,g)
 \le2^{2M(M+1)g+g^2+(M+1)^2g^2}
 \le2^{4P^4}=2^{2^{4\Xi+2}}.
\]
Therefore
\[
 \holeBound(N,d)\le2^{\Xi+d2^{4\Xi+2}}
 \le2^{2^{5\Xi+3}},
\]
because $d\le2^\Xi$ and $\Xi\le2^\Xi\le2^{5\Xi+2}$.

For the specialization, $d_N=8N^2-2N-1$, so $d_N\le8N^2$, $N+d_N\le8N^2$, and
$d_N+1\le8N^2$. Consequently
\[
 4N(N+d_N)(d_N+1)\le256N^5,
 \qquad \Xi(N,d_N)\le256N^5+27N^2+3N+4.
\]
Hence $5\Xi(N,d_N)+3\le1280N^5+135N^2+15N+23\le1280(N+1)^5$, because
$1280(N+1)^5-1280N^5\ge12800N^2+6400N+1280$.
Monotonicity of exponentiation proves \eqref{eq:elementary-closed}.
\end{proof}
